\documentclass[12pt]{article}
\usepackage[utf8]{inputenc}
\usepackage[T1]{fontenc}
\usepackage[margin=1in]{geometry}
\usepackage{setspace}
\usepackage{xcolor}
\usepackage[colorlinks=true, urlcolor=blue!60!black, linkcolor=blue!60!black]{hyperref}
\setlength{\parindent}{0em}
\setlength{\parskip}{0.45em}
\renewcommand{\baselinestretch}{1.0}
\begin{document}
\begin{center}
{\Large \textit{Corfield v. Coryell}, 6 F. Cas. 546 (C.C.E.D. Pa. 1823)\par}
\vspace{0.3em}
{\normalsize Case Brief --- One-Page Edition\par}
\end{center}

\textbf{Memory Jogger.} New Jersey could bar nonresidents from its oyster beds: Article IV protects only ``fundamental'' rights belonging to citizens of all free governments, not a share in a state's common property (\href{https://press-pubs.uchicago.edu/founders/documents/a4_2_1s18.html}{Founders' Constitution}).

\textbf{Facts.} In 1821 armed New Jersey officers seized Corfield's vessel, the \textit{Hiram}, while its sublessee, John Keene, dredged for oysters in Maurice River Cove. A state statute barred nonresidents from gathering oysters in state waters aboard vessels not wholly owned by residents, on pain of forfeiture. Magistrates condemned the vessel, which sold for \$10.

\textbf{Procedural History.} Corfield sued the prize master, Coryell, in trespass in federal circuit court; a \$560 jury verdict was set aside and judgment entered for Coryell.

\textbf{Votes.} A single opinion by Circuit Justice Bushrod Washington.

\textbf{Holding.} The statute violates neither the Commerce, Privileges and Immunities, nor Admiralty Clause, and a lessor out of possession cannot sue in trespass---which makes the constitutional rulings arguably dicta.

\textbf{Reasoning.} Congress's power over navigation did not transfer New Jersey's proprietary interest in its fisheries (\textit{Gibbons v. Ogden}). Article IV secures fundamental rights---travel and residence for trade or professional pursuits, habeas corpus, court access, property, equal taxation, and the franchise as regulated by state law---not every local advantage. A state's citizens hold its fisheries as ``tenants in common,'' and opening them to all would exhaust a finite resource.

\textbf{Examples.} \emph{Same-Side:} higher nonresident elk-hunting fees (\textit{Baldwin}). \emph{Opposite-Side:} a prohibitive fee on nonresident commercial shrimpers (\textit{Toomer}). \emph{Fence-Sitter:} a city's resident-hiring quota on public projects (\textit{Camden}).

\textbf{Critique.} Scholars still divide over whether Washington described substantive natural rights later carried into the Fourteenth Amendment or a mere interstate-equality guarantee, and his common-property fiction can mask protectionism.

\textbf{Key Quotation.} ``Protection by the government; the enjoyment of life and liberty, with the right to acquire and possess property of every kind, and to pursue and obtain happiness and safety; subject nevertheless to such restraints as the government may justly prescribe for the general good of the whole.'' --- Circuit Justice Washington, 6 F. Cas. at 551--52.

\end{document}
